\documentclass{article}
\usepackage{graphicx}
\usepackage[spanish,activeacute,es-tabla]{babel}
\usepackage[utf8x]{inputenc} % mac utf8
\usepackage{amsfonts}
\usepackage{minted}
\usepackage{ucs}

\begin{document}

\title{\textit{Teoría de la Información y Codificación}\\ \textbf{Práctica 2: Creación código con distribución de probabilidad asociada}}
\author{José A. Montenegro Montes}

\maketitle


\section{Enunciado}

La práctica muestra los conceptos teóricos desarrollados en el tema 3. 

A la hora de crear una codificación tenemos que tener en cuenta las propiedades que vimos en el tema anterior. La fuente emisora de los símbolos tiene asociada una distribución de probabilidad que la caracteriza. Dicha distribución de probabilidad debe ser tomada en cuenta a la hora de crear los códigos que realizaremos en esta práctica.

En la práctica además calcularemos un indicador, la entropía, que nos permitirá determina el límite inferior y superior de la longitud del código que estamos buscando.

Para el correcto funcionamiento de la práctica, el alumno deberá implementar los siguientes métodos: 

\begin{itemize}
\item \textit{entropía}. Calcula el valor de la entropía de una distribución dada (página 23 del tema 3).
\item \textit{longitudMedia}. Calcula la longitud media de un código (página 23 del tema 3).
\item \textit{reglaShannonFano}. Obtiene los parámetros $n_i$ para establecer un código posteriormente con el árbol binario (página 39 del tema 3). 
\end{itemize}



\section{Conclusiones}

Con esta práctica observamos como crear códigos adaptados a una fuente generadora. La propiedad que caracteriza a la fuente es la distribución de probabilidad asociada a la emisión de los símbolos. Por tanto, es necesario tener en consideración la distribución de probabilidad a la hora de crear la codificación.

El cálculo de la entropía nos ofrece información sobre el límite inferior en la codificación que podemos obtener. Mediante la regla de Shannon-Fano podemos establecer una codificación ``aceptable'' pero no óptima. Para obtener la codificación óptima es necesario aplicar las dos reglas de Huffman.

La salida del código, una vez implementados los métodos correctamente, será la siguiente:

\begin{minted}[mathescape,
               linenos,
               numbersep=5pt,
               frame=lines,
               framesep=2mm]{csharp}
 
 
Ejercicio 5. Probabilidades: [0.4, 0.2, 0.2, 0.1, 0.1]

Shannon Fano Codigos
--------------------

Code obtained with 0 1 2 2 0  parameters: 
 00  010  011  1000  1001 

 Huffman Codigos
--------------------

Optimal Code obtained with [1, 1, 1, 2, 0, 0, 0, 0, 0, 0] parameters:
[0, 10, 110, 1110, 1111]

****** Calculo Valores ***********

Entropia: 2.1219280948873624
Longitud Media: 2.8
Longitud Optima: 2.2
Verificar h<lo<lm<h+1: true 
 
\end{minted}



\section{Código}


\subsection*{Clase Fuente } 

\inputminted[mathescape,
               numbersep=5pt,
               frame=lines,
               framesep=2mm]{java}{practica2/Fuente.java}

               %%%%%%%%%%%%%%%%%%%%%%%
 
\subsection*{Clase HuffmanCode Java.} 
\inputminted[mathescape,
               numbersep=5pt,
               frame=lines,
               framesep=2mm]{java}{practica2/HuffmanCode.java}
               
%%%%%%



\end{document}
